% =============================================================================
% Moduli-space strata: constructors
% =============================================================================
% Building blocks for the strata of a moduli space of stable curves: dual
% graphs, genus labels, and the poset of strata ordered by closure. The strata
% themselves, their stable graphs and stable curves, are objects, one file
% each, under figures/objects/moduli/, assembled from these constructors; the
% strata of \overline{M}_2 are objects/moduli/m2-*.tikz.
% \input by dzg-constructors.tex, so "@" is a letter here.
%
% Commands:
%   \dualvertex[<label distance>]{<name>}{<x,y>}{<genus>}{<placement>}
%       a vertex of a dual graph with its genus label at <placement>
%       (left, below right, ...) and <label distance> (default 3pt).
%   \dualloop{<x,y>}{left|right}
%       a loop at the vertex at <x,y>: a circle of radius \dualloopradius
%       tangent to the vertex on that side. Draw loops before the vertex so
%       the vertex covers the loop.
%   \genuslabel{<placement>}{<x,y>}{<genus>}
%       the genus of a component of a stable curve, at <x,y>.
%   \strataposet{<object prefix>}{<name/nodes/row, ...>}{<a/b, ...>}
%       the strata of a moduli space as a poset: stratum <name>, with <nodes>
%       nodes (its codimension), at position (nodes, row) of the layout, and
%       a degeneration arrow from a to b for each cover relation a/b
%       (b lies in the closure of a). Stratum <name> is the object
%       <object prefix><render>-<name> placed as -<name>, or with
%       `strata render=name` a poset node labelled <name>; either way its
%       bounding box is -<name>-box, so a stratum object names its bounding
%       box -box. After the call \strataposetdata holds the strata list, for
%       figures that loop over the strata.
%
% Keys:
%   genus labels=all|none       every component and vertex carries its genus,
%                               0 included (default all)
%   strata render=<kind>|name   which object of each stratum \strataposet
%                               places (graph, curve, ...; default graph), or
%                               its name in a poset node
%   strata layout=horizontal|vertical
%                               nodes left to right (x = 4.6 nodes,
%                               y = 1.35 row), or codimension top to bottom
%                               (x = -2.5 row, y = -2 nodes)

\newif\ifdzg@genuslabels
\dzg@genuslabelstrue
\def\dzg@m@render{graph}
\def\dzg@m@layout{horizontal}
\tikzset{
  genus labels/.is choice,
  genus labels/all/.code={\dzg@genuslabelstrue},
  genus labels/none/.code={\dzg@genuslabelsfalse},
  strata render/.store in=\dzg@m@render,
  strata layout/.is choice,
  strata layout/horizontal/.code={\def\dzg@m@layout{horizontal}},
  strata layout/vertical/.code={\def\dzg@m@layout{vertical}},
}

\newcommand\genuslabel[3]{%
  \ifdzg@genuslabels\node[genus, #1] at (#2) {$#3$};\fi}
\newcommand\dualvertex[5][3pt]{%
  \node[dual vertex] (#2) at (#3) {};
  \ifdzg@genuslabels\node[genus, #5=#1] at (#2) {$#4$};\fi}
\def\dualloopradius{0.32}
\def\dzg@m@loop@left{-}
\def\dzg@m@loop@right{}
\newcommand\dualloop[2]{%
  \draw[shift={(#1)}] (\csname dzg@m@loop@#2\endcsname\dualloopradius,0)
    circle (\dualloopradius);}

% Position (\dzg@m@x, \dzg@m@y) of the stratum with <nodes> and <row>.
\def\dzg@m@position#1#2{%
  \csname dzg@m@position@\dzg@m@layout\endcsname{#1}{#2}}
\def\dzg@m@position@horizontal#1#2{%
  \pgfmathsetmacro\dzg@m@x{4.6*#1}\pgfmathsetmacro\dzg@m@y{1.35*#2}}
\def\dzg@m@position@vertical#1#2{%
  \pgfmathsetmacro\dzg@m@x{-2.5*#2}\pgfmathsetmacro\dzg@m@y{-2*#1}}

% \dzg@m@stratum{<object prefix>}{<position>}{<render>}{<name>}
\def\dzg@m@name{name}
\def\dzg@m@stratum#1#2#3#4{%
  \def\dzg@m@r{#3}%
  \ifx\dzg@m@r\dzg@m@name
    \node[poset node] (-#4-box) at (#2) {$#4$};
  \else
    \pic (-#4) at (#2) {object=#1#3-#4};
  \fi}

\newcommand\strataposet[3]{%
  \gdef\strataposetdata{#2}%
  \foreach \dzg@s/\dzg@n/\dzg@row in {#2} {
    \dzg@m@position{\dzg@n}{\dzg@row}
    \edef\dzg@m@next{\noexpand\dzg@m@stratum{#1}{\dzg@m@x,\dzg@m@y}%
      {\dzg@m@render}{\dzg@s}}%
    \dzg@m@next
  }
  \foreach \dzg@a/\dzg@b in {#3}
    \draw[degeneration] (-\dzg@a-box) -- (-\dzg@b-box);
}
